\documentclass[12pt,a4paper]{article}

% --- Layout & Typography ---
\usepackage[top=2.5cm, bottom=2.5cm, left=3cm, right=2.5cm]{geometry}
\usepackage{newtxtext, newtxmath}

% --- Mathematics ---
\usepackage{amsmath}
\usepackage{mathtools}
\usepackage{amssymb}

% --- Units ---
\usepackage{siunitx}
\sisetup{per-mode=symbol, detect-weight=true, detect-family=true}

% --- Tables & Figures ---
\usepackage{booktabs}
\usepackage{tabularx}
\usepackage{graphicx}
\usepackage{float}
\usepackage{subcaption}

% --- Cross-referencing & Bibliography ---
\usepackage[numbers, sort&compress]{natbib}
\usepackage[colorlinks=true, linkcolor=blue, citecolor=blue, urlcolor=blue]{hyperref}
\usepackage[capitalise, noabbrev]{cleveref}

% -------------------------------------------------------------------------
\title{%
  \textbf{Computational Methods for Geotechnical Sieve Analysis}\\[6pt]
  \large Calculation Approaches and Validity Conditions\\
  of the GSD Analysis Pipeline (ASTM D6913)
}
\author{FAFALAB Research Group}
\date{June 2026}
% -------------------------------------------------------------------------

\begin{document}

\maketitle
\tableofcontents
\bigskip\hrule\bigskip

%% =========================================================================
\begin{abstract}
This document describes the computational methods implemented in the
grain-size distribution (GSD) analysis pipeline for geotechnical sieve
testing conducted in accordance with \textbf{ASTM D6913}~\citep{astm_d6913}.
All formulas, interpolation procedures, validity conditions, and screening-level
estimation methods are presented in full.
The pipeline consists of two scripts: \texttt{sieve\_analysis\_test.py}, which
processes raw sieve-mass data and computes geotechnical parameters, and
\texttt{plot\_stratigraphy.py}, which visualises those parameters as a
stratigraphic vertical profile.
Where standard measurement fails due to high fines content (\(>10\%\) passing
the \#200 sieve), log-linear extrapolation and the Kozeny--Carman
equation~\citep{kozeny1927, carman1937} are used to provide indicative
screening estimates that must be confirmed by hydrometer analysis
\citep{astm_d7928}.
\end{abstract}

%% =========================================================================
\section{Introduction}

The primary objective of this document is to provide a transparent, traceable
account of every mathematical operation performed by the GSD pipeline.
It is \emph{not} a site-specific results report; those are generated separately
in \texttt{Site\_N/report/}.

The pipeline follows a two-stage architecture decoupled by JSON:
\begin{enumerate}
    \item \textbf{Stage 1 (\texttt{sieve\_analysis\_test.py})}: Reads raw CSV
          sieve data, computes the cumulative percent-passing curve, interpolates
          characteristic grain diameters, and exports a GSD plot and a JSON
          report per sample.
    \item \textbf{Stage 2 (\texttt{plot\_stratigraphy.py})}: Reads the JSON
          reports for all samples at a site and generates a two-panel
          stratigraphic profile of uniformity coefficient and hydraulic
          conductivity versus depth.
\end{enumerate}

Governing standards referenced throughout this document:
\begin{itemize}
    \item \textbf{ASTM D6913} -- Particle-Size Distribution Using Sieve
          Analysis~\citep{astm_d6913}
    \item \textbf{ASTM D7928} -- Particle-Size Distribution Using Hydrometer
          Analysis~\citep{astm_d7928}
    \item \textbf{ASTM D2487} -- Unified Soil Classification System
          (USCS)~\citep{astm_d2487}
\end{itemize}

%% =========================================================================
\section{Data Input and Pre-Processing}
\label{sec:input}

\subsection{Input Format}

Each sample is supplied as a two-column CSV file with columns
\texttt{Sieve} and \texttt{Sample\_Mass(g)}.  Thirteen sieve designations
are recognised; their physical opening sizes (the values marked on the
laboratory sieve frames) are listed in \cref{tab:sieves}.

\begin{table}[H]
    \centering
    \caption{Sieve designations and physical opening sizes used in the pipeline.}
    \label{tab:sieves}
    \begin{tabular}{lS[table-format=1.3]lS[table-format=1.3]}
        \toprule
        Designation & {Opening (\si{\milli\metre})} &
        Designation & {Opening (\si{\milli\metre})} \\
        \midrule
        \#4   & 4.760 & \#100 & 0.149 \\
        \#10  & 2.000 & \#140 & 0.105 \\
        \#20  & 0.840 & \#200 & 0.074 \\
        \#30  & 0.590 & \#400 & 0.037 \\
        \#40  & 0.420 & Pan   & {---}  \\
        \#50  & 0.297 & & \\
        \#60  & 0.250 & & \\
        \#80  & 0.177 & & \\
        \bottomrule
    \end{tabular}
\end{table}

\noindent
The Pan row acts as a catch-all; its placeholder size (\SI{0.001}{\milli\metre})
is used only for mass accounting and is \emph{excluded} from interpolation
(see \cref{sec:pchip}).

\subsection{Cumulative Percent Passing}
\label{sec:percent_passing}

Following \citet{astm_d6913}, the data are sorted from coarsest to finest sieve.
The cumulative mass retained at each sieve $i$ is
\begin{equation}
    M_{\mathrm{cum},i} = \sum_{j=1}^{i} m_j
    \label{eq:cum_retained}
\end{equation}
where $m_j$ is the mass retained on sieve $j$ (in grams) and the sum proceeds
from the coarsest sieve downward.  The percent passing is then
\begin{equation}
    P_i = 100 \times \left(1 - \frac{M_{\mathrm{cum},i}}{M_{\mathrm{total}}}\right)
    \label{eq:percent_passing}
\end{equation}
where $M_{\mathrm{total}} = \sum_j m_j$ is the total dry specimen mass
including the Pan.  This yields a monotonically decreasing sequence
$P_1 \ge P_2 \ge \cdots \ge P_{13}$ (since the finest sieve always retains
the remaining fines above the Pan).

%% =========================================================================
\section{Grain-Size Interpolation via PCHIP}
\label{sec:pchip}

\subsection{Motivation for Log-Space Interpolation}

Grain-size distribution curves are conventionally plotted on a semi-logarithmic
axis (particle size on a log scale, percent passing on a linear scale).
Interpolation performed in log-space therefore respects the shape of the
curve and avoids the artefacts that linear-space interpolation introduces at
the coarser end of the distribution.

\subsection{PCHIP Method}

The pipeline uses the \emph{Piecewise Cubic Hermite Interpolating Polynomial}
(PCHIP) algorithm~\citep{fritsch1980}, which guarantees monotonicity of the
interpolant between successive data points.  This is an important property
for grain-size curves: spurious oscillations (overshoot) must be avoided
because percent-passing values must remain in \([0, 100]\).

Formally, let $\{(P_i,\, d_i)\}$ be the set of measured (percent passing,
grain diameter) pairs with the Pan excluded.  PCHIP is applied to the
log-transformed diameters:
\begin{equation}
    D_x = 10^{\,\mathrm{PCHIP}\!\left(x,\;\log_{10}(d)\right)}
    \label{eq:pchip}
\end{equation}
where $x$ is the target percentile (e.g.\ $x = 10$ for $D_{10}$) and
$\mathrm{PCHIP}(\cdot)$ denotes the monotone cubic Hermite interpolant fitted
to the pairs $\{(P_i,\, \log_{10} d_i)\}$.

\subsection{Validity Conditions}

Interpolation is strictly limited to the measured data range.
If the target percentile $x$ does not satisfy
\begin{equation}
    \min_i(P_i) \;\le\; x \;\le\; \max_i(P_i),
    \label{eq:validity_range}
\end{equation}
the function raises a \texttt{ValueError} and the corresponding $D_x$ value is
reported as \texttt{null} in the JSON output.  No extrapolation is performed
within the measured range; when the curve cannot reach a percentile, a
hydrometer analysis per \citet{astm_d7928} is required to extend the GSD
curve into the fine-grained (silt and clay) domain.

\subsection{Target Percentiles}

The pipeline computes $D_{10}$, $D_{25}$, $D_{30}$, $D_{50}$, $D_{60}$, and
$D_{75}$.  $D_{50}$ is used internally for the Kozeny--Carman estimate
(\cref{sec:kc}) but is not reported in the public JSON output to avoid
confusion with the directly computed characteristic diameters.

%% =========================================================================
\section{Measured Geotechnical Coefficients}
\label{sec:measured}

The following coefficients are computed from the PCHIP-interpolated $D_x$
values.  Any coefficient whose prerequisite $D_x$ is \texttt{null}
(i.e.\ outside the measurable range) is itself reported as \texttt{null}.

\subsection{Uniformity Coefficient \texorpdfstring{$C_u$}{Cu}}

\begin{equation}
    C_u = \frac{D_{60}}{D_{10}}
    \label{eq:cu}
\end{equation}

$C_u$ quantifies the spread of the grain-size distribution.  A soil with
$C_u = 1$ is perfectly uniform (single grain size); larger values indicate
a broader, better-graded distribution.  Per \citet{astm_d2487}:
\begin{itemize}
    \item \emph{Well-graded sand} (SW): $C_u \ge 6$ and $1 \le C_c \le 3$
    \item \emph{Well-graded gravel} (GW): $C_u \ge 4$ and $1 \le C_c \le 3$
\end{itemize}

\subsection{Coefficient of Curvature \texorpdfstring{$C_c$}{Cc}}

\begin{equation}
    C_c = \frac{(D_{30})^2}{D_{60} \cdot D_{10}}
    \label{eq:cc}
\end{equation}

$C_c$ describes the shape of the central portion of the GSD curve.  The
USCS well-graded criterion requires $1 \le C_c \le 3$ in conjunction with
the $C_u$ threshold stated above \citep{astm_d2487}.

\subsection{Trask Sorting Coefficient \texorpdfstring{$S_0$}{S0}}

\begin{equation}
    S_0 = \sqrt{\frac{D_{75}}{D_{25}}}
    \label{eq:s0}
\end{equation}

The Trask sorting coefficient~\citep{trask1932} is a statistical measure of
grain-size dispersion based on the geometric quartile deviation.  Interpretation
guidelines are summarised in \cref{tab:s0}.

\begin{table}[H]
    \centering
    \caption{Trask sorting coefficient interpretation \citep{trask1932}.}
    \label{tab:s0}
    \begin{tabular}{ll}
        \toprule
        $S_0$ range & Classification \\
        \midrule
        $< 2.5$     & Well-sorted (uniform grains) \\
        $\approx 3$ & Normally sorted \\
        $> 4.5$     & Poorly sorted (wide grain-size range) \\
        \bottomrule
    \end{tabular}
\end{table}

\subsection{Hazen Hydraulic Conductivity \texorpdfstring{$K_{\mathrm{Hazen}}$}{K\_Hazen}}
\label{sec:hazen}

The empirical Hazen formula~\citep{hazen1911} relates hydraulic conductivity
to the effective size $D_{10}$:
\begin{equation}
    K_{\mathrm{Hazen}} = (D_{10})^2 \quad [\si{\centi\metre\per\second}]
    \label{eq:hazen}
\end{equation}
where the Hazen constant $C = 1.0$ is assumed (appropriate for medium, clean
sands).  The formula is applied only when \emph{all} of the following
conditions are satisfied:
\begin{enumerate}
    \item $D_{10}$ is within the measurable sieve range (not \texttt{null}),
    \item $\SI{0.1}{\milli\metre} \le D_{10} \le \SI{3.0}{\milli\metre}$
          (effective-size validity window),
    \item $C_u < 5$ (uniform-sand criterion for the Hazen approximation).
\end{enumerate}
If any condition is unmet, $K_{\mathrm{Hazen}}$ is reported as \texttt{null}
and the Kozeny--Carman estimate in \cref{sec:kc} serves as a supplementary
indicator.

%% =========================================================================
\section{Estimated Parameters for High-Fines Samples}
\label{sec:estimated}

\subsection{Problem Statement}

When more than approximately 10\% of the sample passes the \#200 sieve
(\SI{0.074}{\milli\metre}), the effective grain diameter $D_{10}$ falls below
the measurable range of dry sieve analysis.  Consequently, $C_u$, $C_c$, and
$K_{\mathrm{Hazen}}$ cannot be computed from measurements alone.

ASTM D2487 acknowledges this limitation and permits extrapolation for the
purpose of obtaining $D_{10}$~\citep{astm_d2487}; however, any extrapolated
value is indicative only and must be confirmed by hydrometer analysis
\citep{astm_d7928}.

\subsection{Log-Linear \texorpdfstring{$D_{10}$}{D10} Extrapolation}
\label{sec:d10est}

When $D_{10}$ is \texttt{null}, the pipeline extrapolates using the three sieve
points with the lowest percent-passing values (i.e.\ the three finest physical
sieves with measurable mass).  In log-diameter versus percent-passing space,
these points are approximated by a straight line:
\begin{align}
    \log_{10}(d_j) &= m \cdot p_j + b, \quad j = 1, 2, 3
    \label{eq:linreg} \\
    D_{10,\mathrm{est}} &= 10^{\,m \cdot 10 + b}
    \label{eq:d10est}
\end{align}
where $(m, b)$ are the least-squares slope and intercept obtained from
\texttt{numpy.polyfit} applied to the three (percent-passing, log-diameter)
pairs.

A validity guard is applied: the extrapolated value is accepted only if
\begin{equation}
    \SI{0.001}{\milli\metre} < D_{10,\mathrm{est}} < \SI{0.074}{\milli\metre}
    \label{eq:d10est_validity}
\end{equation}
ensuring the estimate lies within the silt range (below the \#200 sieve but
above the Pan placeholder size).  If the guard fails, $D_{10,\mathrm{est}}$
is reported as \texttt{null}.

\subsection{Estimated \texorpdfstring{$C_{u,\mathrm{est}}$}{Cu\_est} and
            \texorpdfstring{$C_{c,\mathrm{est}}$}{Cc\_est}}

Given a valid $D_{10,\mathrm{est}}$, the measured $D_{30}$ and $D_{60}$
(which are typically available for samples with moderate-to-high fines) are
substituted into the standard formulas:
\begin{equation}
    C_{u,\mathrm{est}} = \frac{D_{60}}{D_{10,\mathrm{est}}}
    \label{eq:cu_est}
\end{equation}
\begin{equation}
    C_{c,\mathrm{est}} = \frac{(D_{30})^2}{D_{60} \cdot D_{10,\mathrm{est}}}
    \label{eq:cc_est}
\end{equation}
These values are labelled \enquote*{Estimated (indicative only)} in both the
JSON output and the GSD plot, and carry a warning directing the user to
confirm with ASTM D7928.

\subsection{Kozeny--Carman Hydraulic Conductivity \texorpdfstring{$K_{\mathrm{KC}}$}{K\_KC}}
\label{sec:kc}

For samples where $D_{50}$ is within the measurable range (i.e.\ less than
50\% fines), the Kozeny--Carman equation~\citep{kozeny1927, carman1937}
provides a physics-based estimate of hydraulic conductivity:
\begin{equation}
    K_{\mathrm{KC}}(n) =
    \frac{g}{\nu}
    \cdot \frac{D_{50,\mathrm{m}}^{2}}{180}
    \cdot \frac{n^3}{(1-n)^2}
    \times 100
    \quad [\si{\centi\metre\per\second}]
    \label{eq:kc}
\end{equation}
where:
\begin{itemize}
    \item $g = \SI{9.81}{\metre\per\second\squared}$ -- gravitational acceleration,
    \item $\nu = \SI{1e-6}{\metre\squared\per\second}$ -- kinematic viscosity of
          water at \SI{20}{\degreeCelsius},
    \item $D_{50,\mathrm{m}} = D_{50} \times 10^{-3}$ -- median grain diameter
          converted from \si{\milli\metre} to \si{\metre},
    \item $n$ -- porosity (dimensionless).
\end{itemize}

Natural sandy soils typically exhibit porosities in the range
$n \in [0.25, 0.50]$~\citep{astm_d2487}.  The pipeline uses the narrower
range $n \in [0.40, 0.45]$ representative of loose-to-medium dense sands,
and reports two bounding estimates:
\begin{align}
    K_{\mathrm{KC,low}}  &= K_{\mathrm{KC}}(n = 0.40)
    \label{eq:kc_low} \\
    K_{\mathrm{KC,high}} &= K_{\mathrm{KC}}(n = 0.45)
    \label{eq:kc_high}
\end{align}
This bracketing approach communicates the sensitivity of the estimate to
porosity uncertainty without requiring a direct measurement of void ratio.

%% =========================================================================
\section{Stratigraphic Profile Computation}
\label{sec:stratigraphy}

The script \texttt{plot\_stratigraphy.py} reads all Site~1 JSON reports and
assembles a two-panel stratigraphic profile.

\subsection{Depth Convention}

Sample names follow the convention \texttt{Sample\_1-\textit{H}}, where
$H$ is the height in metres above the base of the \SI{10}{\metre} excavation
pit.  The depth is extracted by parsing the integer suffix:
\[
    H = \mathrm{int}\!\left(\texttt{sample\_name.split(`-')[-1]}\right)
\]
Thus $H = 0$ corresponds to the deepest (basal) layer and $H = 10$ to the
surface layer.

\subsection{Most-Permeable Layer Identification}

To highlight the stratigraphic horizon of greatest hydraulic interest, the
script identifies the layer with the maximum measured $K_{\mathrm{Hazen}}$:
\begin{equation}
    i^{*} = \operatorname*{arg\,max}_{i} \; K_{\mathrm{Hazen},i}
    \label{eq:anomaly}
\end{equation}
This layer is annotated on both panels of the profile as
\enquote*{Most Permeable Layer}.  Note that samples with \texttt{null}
$K_{\mathrm{Hazen}}$ (high-fines samples) are excluded from the
$\operatorname{arg\,max}$ operation via \texttt{numpy.nanargmax}.

\subsection{Plot Layout}

The output figure contains two subplots sharing the same $y$-axis (height
above pit base):
\begin{itemize}
    \item \textbf{Left panel}: $C_u$ on a linear scale.
    \item \textbf{Right panel}: $K_{\mathrm{Hazen}}$ on a base-10 logarithmic
          scale (\si{\centi\metre\per\second}).
\end{itemize}
The $y$-axis is not inverted; $H = 0$ is at the bottom and $H = 10$ is at
the top, matching the physical orientation of the pit.

%% =========================================================================
\section{Summary of Formulas and Validity Conditions}
\label{sec:summary}

\cref{tab:summary} consolidates all formulas, their equation numbers, and the
conditions under which each quantity is computable.

\begin{table}[H]
    \centering
    \caption{Summary of all formulas, references, and validity conditions.}
    \label{tab:summary}
    \renewcommand{\arraystretch}{1.35}
    \begin{tabularx}{\linewidth}{l l l X}
        \toprule
        Parameter & Formula & Ref. & Validity / Notes \\
        \midrule
        $P_i$ (percent passing)
            & \eqref{eq:percent_passing}
            & \citep{astm_d6913}
            & Always computable \\
        $D_x$ (grain diameter)
            & \eqref{eq:pchip}
            & \citep{fritsch1980}
            & $\min(P) \le x \le \max(P)$; Pan excluded \\
        $C_u$ (uniformity)
            & \eqref{eq:cu}
            & \citep{astm_d2487}
            & Requires $D_{10}$, $D_{60}$ \\
        $C_c$ (curvature)
            & \eqref{eq:cc}
            & \citep{astm_d2487}
            & Requires $D_{10}$, $D_{30}$, $D_{60}$ \\
        $S_0$ (Trask sorting)
            & \eqref{eq:s0}
            & \citep{trask1932}
            & Requires $D_{25}$, $D_{75}$ \\
        $K_{\mathrm{Hazen}}$
            & \eqref{eq:hazen}
            & \citep{hazen1911}
            & $D_{10}$ available; \SIrange{0.1}{3.0}{\milli\metre}; $C_u < 5$ \\
        $D_{10,\mathrm{est}}$
            & \eqref{eq:d10est}
            & \citep{astm_d2487}
            & $D_{10}$ = null; result in \SIrange{0.001}{0.074}{\milli\metre} \\
        $C_{u,\mathrm{est}}$
            & \eqref{eq:cu_est}
            & ---
            & $D_{10,\mathrm{est}}$ valid; confirm with \citet{astm_d7928} \\
        $C_{c,\mathrm{est}}$
            & \eqref{eq:cc_est}
            & ---
            & $D_{10,\mathrm{est}}$ valid; confirm with \citet{astm_d7928} \\
        $K_{\mathrm{KC}}$ (low/high)
            & \eqref{eq:kc}
            & \citep{kozeny1927, carman1937}
            & $D_{50}$ available; $n = 0.40$--$0.45$; indicative only \\
        $i^{*}$ (most permeable)
            & \eqref{eq:anomaly}
            & ---
            & NaN-safe \texttt{nanargmax} over $K_{\mathrm{Hazen}}$ \\
        \bottomrule
    \end{tabularx}
\end{table}

%% =========================================================================
\bibliographystyle{plainnat}
\bibliography{references}

\end{document}
